\chapter{Ablations}
\label{sec:ablations}
We ablate the main design choices of ECASP on the MVPSP test set using the baseline checkpoint, the same synthetic-training protocol, and benchmark-provided real masks as in Sec.~\ref{sec:main_results}.
All errors are in millimeters except \(\Delta R\), which is in degrees.

\noindent\textbf{Number of epipolar samples.}
Tab.~\ref{tab:ablation_num_epipolar_samples} shows that \(S=24\) provides the best overall balance.
Increasing to \(S=32\) reduces translation error but degrades rotation and ADD/ADD-S, indicating that denser epipolar sampling does not necessarily yield better pose estimates.

\begin{table}[htbp]
\centering
\caption{Ablation on the number of sampled points \(S\) along each epipolar line.}
\label{tab:ablation_num_epipolar_samples}
\small
\setlength{\tabcolsep}{4.5pt}
\begin{tabular*}{\linewidth}{@{\extracolsep{\fill}}ccccccc@{}}
\toprule
\multirow{2}{*}{\(S\)}
&
\multicolumn{3}{c}{Drill}
&
\multicolumn{3}{c}{Screwdriver}
\\
\cmidrule(lr){2-4}
\cmidrule(lr){5-7}
&
ADD \(\downarrow\)
&
\(\Delta t \downarrow\)
&
\(\Delta R \downarrow\)
&
ADD-S \(\downarrow\)
&
\(\Delta t \downarrow\)
&
\(\Delta R \downarrow\)
\\
\midrule
\(16\)
& \(5.62\)
& \(1.24\)
& \(3.34\)
& \(2.78\)
& \(2.67\)
& \(4.27\)
\\
\(24\)
& \(\mathbf{4.93}\)
& \(1.33\)
& \(\mathbf{2.88}\)
& \(\mathbf{2.47}\)
& \(2.02\)
& \(\mathbf{3.76}\)
\\
\(32\)
& \(6.05\)
& \(\mathbf{1.23}\)
& \(3.57\)
& \(2.64\)
& \(\mathbf{1.89}\)
& \(4.07\)
\\
\bottomrule
\end{tabular*}
\end{table}

\noindent\textbf{Auxiliary losses.}
Tab.~\ref{tab:ablation_losses} shows that the auxiliary losses consistently improve the asymmetric drill, particularly ADD and rotation.
The effect is smaller for the symmetric screwdriver, where some rotational differences are less observable under the evaluation symmetry.
This suggests that the auxiliary objectives provide stronger additional constraints when the object pose is less ambiguous.


\begin{table}[htbp]
\centering
\caption{Ablation on the auxiliary training losses.}
\label{tab:ablation_losses}
\small
\setlength{\tabcolsep}{4.5pt}
\begin{tabular*}{\linewidth}{@{\extracolsep{\fill}}lcccccc@{}}
\toprule
\multirow{2}{*}{Variant}
&
\multicolumn{3}{c}{Drill}
&
\multicolumn{3}{c}{Screwdriver}
\\
\cmidrule(lr){2-4}
\cmidrule(lr){5-7}
&
ADD \(\downarrow\)
&
\(\Delta t \downarrow\)
&
\(\Delta R \downarrow\)
&
ADD-S \(\downarrow\)
&
\(\Delta t \downarrow\)
&
\(\Delta R \downarrow\)
\\
\midrule
Full
& \(\mathbf{4.93}\)
& \(\mathbf{1.33}\)
& \(\mathbf{2.88}\)
& \(2.47\)
& \(2.02\)
& \(\mathbf{3.76}\)
\\
w/o \(\mathcal{L}_{\mathrm{mc}}\)
& \(5.38\)
& \(1.36\)
& \(3.19\)
& \(2.89\)
& \(2.75\)
& \(3.99\)
\\
w/o \(\mathcal{L}_{\mathrm{allo}}\)
& \(5.50\)
& \(1.37\)
& \(3.24\)
& \(2.61\)
& \(2.00\)
& \(3.82\)
\\
w/o both
& \(5.55\)
& \(1.36\)
& \(3.30\)
& \(\mathbf{2.43}\)
& \(\mathbf{1.99}\)
& \(3.90\)
\\
\bottomrule
\end{tabular*}
\end{table}


\noindent\textbf{Epipolar cross-view attention.}
Tab.~\ref{tab:ablation_epipolar_attention} shows that epipolar cross-view attention
outperforms unconstrained cross-attention on all metrics, with the largest gains
in translation.
% This supports the role of calibrated epipolar constraints in cross-view fusion.
This demonstrates the benefit of calibrated epipolar constraints in guiding multi-view feature interaction.

\begin{table}[htbp]
\centering
\caption{Ablation on the epipolar cross-view attention mechanism.}
\label{tab:ablation_epipolar_attention}
\small
\setlength{\tabcolsep}{4.5pt}
\begin{tabular*}{\linewidth}{@{\extracolsep{\fill}}lcccccc@{}}
\toprule
\multirow{2}{*}{Variant}
&
\multicolumn{3}{c}{Drill}
&
\multicolumn{3}{c}{Screwdriver}
\\
\cmidrule(lr){2-4}
\cmidrule(lr){5-7}
&
ADD \(\downarrow\)
&
\(\Delta t \downarrow\)
&
\(\Delta R \downarrow\)
&
ADD-S \(\downarrow\)
&
\(\Delta t \downarrow\)
&
\(\Delta R \downarrow\)
\\
\midrule
Vanilla
& \(6.59\)
& \(2.50\)
& \(3.54\)
& \(2.87\)
& \(3.14\)
& \(4.11\)
\\
Epipolar
& \(\mathbf{4.93}\)
& \(\mathbf{1.33}\)
& \(\mathbf{2.88}\)
& \(\mathbf{2.47}\)
& \(\mathbf{2.02}\)
& \(\mathbf{3.76}\)
\\
\bottomrule
\end{tabular*}
\end{table}


\noindent\textbf{Number of input views.}
Tab.~\ref{tab:ablation_num_views} compares four views with the two-view OL+OR setting.
Four views improve the drill across all metrics, whereas two views perform
slightly better on the symmetric screwdriver, likely due to less informative
or occluded additional views.
We retain four views for their clear gains on the drill and comparability with
prior MVPSP multi-view evaluations~\cite{haugaard2023episurfemb}.


\begin{table}[htbp]
\centering
\caption{Ablation on the number of input views.}
\label{tab:ablation_num_views}
\small
\setlength{\tabcolsep}{4.5pt}
\begin{tabular*}{\linewidth}{@{\extracolsep{\fill}}lcccccc@{}}
\toprule
\multirow{2}{*}{Views}
&
\multicolumn{3}{c}{Drill}
&
\multicolumn{3}{c}{Screwdriver}
\\
\cmidrule(lr){2-4}
\cmidrule(lr){5-7}
&
ADD \(\downarrow\)
&
\(\Delta t \downarrow\)
&
\(\Delta R \downarrow\)
&
ADD-S \(\downarrow\)
&
\(\Delta t \downarrow\)
&
\(\Delta R \downarrow\)
\\
\midrule
OL+OR
& \(6.29\)
& \(1.46\)
& \(3.53\)
& \(\mathbf{2.31}\)
& \(\mathbf{1.64}\)
& \(\mathbf{3.33}\)
\\
L+OL+OR+R
& \(\mathbf{4.93}\)
& \(\mathbf{1.33}\)
& \(\mathbf{2.88}\)
& \(2.47\)
& \(2.02\)
& \(3.76\)
\\
\bottomrule
\end{tabular*}
\end{table}